% !TEX root = ../../hw03.tex
% Homework 03, question 6.
\begin{proof}[Solution]
  We use $\R_+=[0,\infty)$ and positive rational exponents. The
  nonnegative $n$th root exists and is unique for every positive
  integer $n$.

  For part (a), put $u=x^{1/n}$ and $v=y^{1/n}$. Both are nonnegative,
  and the integer power laws give
  \[
    (uv)^n=u^nv^n=xy.
  \]
  Thus $uv$ is the nonnegative $n$th root of $xy$. Uniqueness gives
  \[
    (xy)^{1/n}=x^{1/n}y^{1/n}.
  \]

  For part (b), put $w=(x^{1/n})^m$. This number is nonnegative, and
  \[
    w^n=(x^{1/n})^{mn}=\bigl((x^{1/n})^n\bigr)^m=x^m.
  \]
  Hence $w$ is the nonnegative $n$th root of $x^m$. Again by
  uniqueness,
  \[
    x^{m/n}=(x^m)^{1/n}=(x^{1/n})^m.
  \]

  For part (c), write $r=m/n$ and $s=p/q$, where $m,n,p,q$ are
  positive integers. We use the fact that rational powers are
  independent of the fraction used to represent the exponent.
  First let $u=(x^r)^s$. By definition, $u\geq0$ and
  $u^q=(x^r)^p$. Taking integer powers gives
  \[
    u^{nq}=\bigl((x^r)^p\bigr)^n
      =\bigl((x^r)^n\bigr)^p=x^{mp}.
  \]
  Since $rs=mp/(nq)$, the number $x^{rs}$ is also the nonnegative
  $(nq)$th root of $x^{mp}$. Therefore
  \[
    (x^r)^s=x^{rs}.
  \]

  For the sum law, let $v=x^rx^s$. This is nonnegative, and
  \begin{align*}
    v^{nq}
      &=(x^r)^{nq}(x^s)^{nq} \\
      &=\bigl((x^r)^n\bigr)^q\bigl((x^s)^q\bigr)^n \\
      &=x^{mq}x^{pn}=x^{mq+pn}.
  \end{align*}
  But $r+s=(mq+pn)/(nq)$, so $x^{r+s}$ is the nonnegative $(nq)$th
  root of the same number. Uniqueness therefore gives
  \[
    x^{r+s}=x^rx^s.
  \]
  These arguments include $x=0$, since all exponents are positive
  and we never divide by $x$.
\end{proof}
